% Supplementary practice source only; outside the 100-mark core paper.
% Requires the shared midterm style when included in another document.
\newpage
\textbf{\Large Supplementary practice}\par
S1-S3 cover question types not selected for the main paper. They are outside its 100 marks and 3-hour time limit. Questions 1-12 remain the timed mock examination.\par
\medskip
\textbf{S1\quad Use Case Diagram: University Grades Submission}\par
\textbf{Requirements} (adapted from Lecture 4, pages 64-65)\par
\begin{enumerate}[itemsep=3pt]
\item Instructors log in with their KU account through KU Single Sign-On. Only valid accounts may access the system.
\item Instructors choose a semester and view the courses and sections offered in their department, supplied by the KU Registration System.
\item For a selected course and section(s), the system can load the registered students from the KU Registration System into a grade sheet.
\item A grade sheet has a header (Course, Semester) and student rows. It starts in Draft. Instructors enter or update grades and may print a copy for their records.
\item Once all grades are entered, the Instructor changes the sheet to Ready. Submission is allowed only in Ready, not Draft.
\item On submission, the system updates every student's transcript, emails each student their grade, confirms success, and changes the sheet to Submitted.
\end{enumerate}
\textbf{The five external systems}\par
\begin{tabularx}{\linewidth}{|>{\raggedright\arraybackslash}m{5.3cm}|Y|}\hline
KU Single Sign-On & Verifies accounts and manages passwords\\\hline
KU Registration System & Supplies courses, sections and registered students\\\hline
KU Transcript System & Records official grades\\\hline
KU Student Information System & Supplies contact details, such as student email addresses\\\hline
KU Email Service (Google Workspace) & Sends grade notifications\\\hline
\end{tabularx}
\par\medskip
\textbf{Task:} Draw a Use Case Diagram on the next page covering the Instructor's goals and the external systems. Show the system boundary, actors and associations. Where justified by the requirements, show \code{<<include>>}, \code{<<extend>>}, or generalization with the correct direction and notation. State the condition for any \code{<<extend>>}; do not invent relationships unsupported by the scenario or replace use cases with workflow arrows.\par
\textbf{Brief explanation:} Explain whether each relationship shown is mandatory or optional, what permits submission, and whether the requirements support any generalization.\par
\blankbox{3.5cm}
\newpage
\textbf{\Large S1 (continued): Use Case Diagram workspace}\par
University Grades Submission\hfill Supplementary; outside the 100 marks\par
\blankbox{21.5cm}
\newpage
\textbf{\Large S2\quad State Diagram: Service Ticket}\par
Supplementary; outside the 100 marks. Use the worked path in Lecture 4, pages 61-62, with the supporting facts on pages 45-46.\par
\begin{enumerate}[itemsep=2pt]
\item When a new Ticket is saved, the System automatically creates its first Action Record, sets the status to Requested, and copies Initial Due Date to New Due Date.
\item The CCE may act only on Requested tickets: approve to Approved or close to Closed by CCE, with notes or a new due date where applicable.
\item In the worked path, a Programmer starts an Approved ticket, making it In Progress. Further progress actions may record notes or a new due date while keeping In Progress.
\item When work is finished, the Programmer closes it to Closed by Programmer and sets New Due Date to the completion date.
\item Every saved action updates the Ticket's Last Status and Last Due Date from that latest action. Saved actions cannot be edited or deleted.
\item No new actions are allowed after Closed by CCE or Closed by Programmer. Customers and IT Admins do not perform these status-changing actions.
\end{enumerate}
\textbf{Task:} Draw only the status lifecycle of one Ticket. Show initial/final nodes, states, transitions and a self-loop. Label each transition with its event and the actor in brackets, following the sample notation, e.g. \code{approve [CCE]}. Customer, CCE and Programmer are not states. Do not invent additional states.\par
\blankbox{13cm}
\newpage
\textbf{\Large S3\quad Git: Explain the repository graph}\par
Supplementary; outside the 100 marks. Adapted from Samples Question 2.2, pages 8-9. Use ExamScore and the Git assumptions from Questions 11-12. Read each question independently from this snapshot unless a command is explicitly run.\par
\begin{center}\supplementgitgraph\end{center}
\textbf{S3.1} Which local branch is checked out, and which commit does HEAD identify?\par
\blankbox{1.5cm}
\textbf{S3.2} Which commit at the local sprint tip is not in the history referenced by origin/sprint\_1\_staging? Explain using the pointers.\par
\blankbox{1.8cm}
\textbf{S3.3} Which commit supplies the files in the current working directory?\par
\blankbox{1.5cm}
\textbf{S3.4} Can \code{git merge feature/load\_scores} fast-forward now? What merge would be needed, and why?\par
\blankbox{2cm}
\textbf{S3.5} Does the graph prove that this merge would cause a file conflict? Explain.\par
\blankbox{1.5cm}
\textbf{S3.6} What may \code{git fetch origin} update, and what will it not move automatically?\par
\blankbox{2cm}
